\chapter{Repeating Tasks in Python: Loops}

\section{Introduction: Real-Life Hook}

Students repeat tasks every day:  

\begin{itemize}
    \item Brushing teeth every morning  
    \item Practicing a musical instrument  
    \item Counting items in a collection  
\end{itemize}

Python can repeat tasks automatically using \textbf{loops}, making programs shorter and more efficient.

\section{Concept: Loops in Python}

\begin{definition}{Loop}
A \textbf{loop} is a sequence of instructions that repeats until a condition is met or for a specific number of times.
\end{definition}

\subsection{For Loop: Repeating a Fixed Number of Times}

\begin{verbatim}
# Print numbers 1 to 5
for i in range(1, 6):
    print(i)
\end{verbatim}

Output:

\begin{verbatim}
1
2
3
4
5
\end{verbatim}

- \texttt{range(1,6)} generates numbers from 1 to 5.  
- \texttt{i} is the loop variable that changes in each iteration.

\subsection{While Loop: Repeating Until a Condition is True}

\begin{verbatim}
# Count from 1 to 5
count = 1
while count <= 5:
    print(count)
    count = count + 1
\end{verbatim}

Output is the same as the for loop.

\section{Real-Life Analogy: Morning Routine}

\begin{itemize}
    \item Brush teeth $\rightarrow$ repeat for 2 minutes  
    \item Jumping jacks $\rightarrow$ repeat 10 times  
    \item Practice piano $\rightarrow$ repeat scales until 5 minutes pass  
\end{itemize}

\section{Step-by-Step Diagram: Loop Execution}

\begin{center}
\begin{tikzpicture}[node distance=1.5cm]
\node (start) [draw, ellipse, fill=green!20] {Start Loop};
\node (condition) [draw, diamond, fill=yellow!20, below of=start] {Condition True?};
\node (task) [draw, rectangle, fill=orange!20, below of=condition, yshift=-1cm] {Perform Task};
\node (update) [draw, rectangle, fill=orange!20, right of=task, xshift=3cm] {Update Counter/Variable};
\node (end) [draw, ellipse, fill=green!20, below of=task, yshift=-1.5cm] {End Loop};
\draw[->] (start) -- (condition);
\draw[->] (condition.south) -- node[midway, right] {Yes} (task.north);
\draw[->] (task.east) -- (update.west);
\draw[->] (update.north) |- (condition.east);
\draw[->] (condition.west) -- ++(-1.5,0) |- node[near start, left] {No} (end.west);
\end{tikzpicture}
\end{center}

\section{Mini-Project: Multiplication Table}

Create a Python program that prints the multiplication table of a number:

\begin{verbatim}
num = int(input("Enter a number: "))
for i in range(1, 11):
    print(num, "x", i, "=", num * i)
\end{verbatim}

\section{Using Loops with Lists}

Loops are especially useful when working with lists:

\begin{verbatim}
fruits = ["apple", "banana", "mango"]
for fruit in fruits:
    print("I like", fruit)
\end{verbatim}

Output:

\begin{verbatim}
I like apple
I like banana
I like mango
\end{verbatim}

\section{Key Concepts}

\begin{itemize}
    \item Loops allow repetition of tasks in programs.  
    \item Use \texttt{for} for a known number of iterations, \texttt{while} for condition-based repetition.  
    \item Loop variables control iteration.  
    \item Loops can be combined with lists for efficient processing.  
\end{itemize}

\section{End-of-Chapter Exercises}

\begin{exercise}{Very Short Questions}
\begin{enumerate}
    \item What is a loop in Python?  
    \item Which loop is used for a fixed number of repetitions?  
    \item Name the Python function to generate a sequence of numbers for a for loop.  
\end{enumerate}
\end{exercise}

\begin{exercise}{Short Answer Questions}
\begin{enumerate}
    \item Write a for loop that prints all even numbers from 2 to 10.  
    \item Write a while loop that counts down from 5 to 1.  
    \item Explain the difference between for and while loops.  
\end{enumerate}
\end{exercise}

\begin{exercise}{Activity-Based Questions}
\begin{enumerate}
    \item Create a program that asks the user for 5 numbers and prints their sum using a loop.  
    \item Use a loop to print all items in your favorite snack list.  
    \item Modify the multiplication table program to print tables from 1 to 5.  
\end{enumerate}
\end{exercise}

\section{Teacher Notes}

\marginnote{
\textbf{Learning Objectives:}  
\begin{itemize}
    \item Students understand loops and their purpose.  
    \item Students learn to use for and while loops in Python.  
    \item Students can iterate over numbers and lists.  
\end{itemize}

\textbf{Teaching Tips:}  
\begin{itemize}
    \item Demonstrate both for and while loops.  
    \item Encourage students to try modifying ranges and loop variables.  
\end{itemize}

\textbf{Assessment:}  
\begin{itemize}
    \item Check correctness of loops and outputs.  
    \item Observe understanding of loop repetition and variable updates.  
\end{itemize}
}